Write $CD=d$, $AC=e$, $BD=f$, and let $\alpha=\angle BAD$. Opposite angles of a cyclic quadrilateral are supplementary. Applying the cosine law in $BAD$ and $BCD$ gives
\[
f^2=29-20\cos\alpha=25+d^2+10d\cos\alpha.
\]
Consequently $\cos\alpha=\dfrac{(2-d)}{10}$ and $f^2=25+2d$. The good-quadrilateral condition now gives $e^2=d^2-2d$, so $d>2$. Ptolemy's theorem says $ef=5(d+2)$. Squaring and substituting yields
\[
d(d-2)(25+2d)=25(d+2)^2,
\qquad (d-10)(d^2+8d+5)=0.
\]
Only $d=10$ is positive and greater than $2$. Thus $AC=4\sqrt5$. Triangle $ABC$ has equal sides $5$ and altitude $\sqrt5$ to $AC$, so its area is $10$. Its circumradius is
\[
R=\dfrac{5\cdot5\cdot4\sqrt5}{4\cdot10}=\dfrac{5\sqrt5}{2}.
\]
The extended sine rule in $ACD$ gives $\sin^2\angle ACD=\dfrac{AD^2}{(4R^2)}=\dfrac{4}{125}$. Hence $\boxed{m+n=129}$, choice E.
\begin{asy}
texpreamble("\usepackage{amsmath,amssymb}");
size(240);
defaultpen(fontsize(9pt));
pair A=(-2sqrt(5),sqrt(5)),B=(0,0),C=(2sqrt(5),sqrt(5)),D=(-12sqrt(5)/5,9sqrt(5)/5),O=(0,5sqrt(5)/2);draw(circle(O,5sqrt(5)/2),gray);draw(A--B--C--D--cycle);draw(A--C,blue);draw(B--D,dashed);label("$A$",A,(-1,0));label("$B$",B,(0,-1));label("$C$",C,(1,0));label("$D$",D,NW);
\end{asy}
